\section{Open Questions for Follow-Up}

Five substantive, unresolved questions raised by (or left un-answered by) this
replication. None of these are ``run the same code again with more $t$-points'';
each requires new implementation or a new comparison basis.

\begin{enumerate}
  \item \textbf{Does the MPF-vs-Suzuki-4 crossover advantage persist quantitatively for
        physically-motivated Hamiltonians (TFIM, Fermi--Hubbard, small molecular electronic
        structure) at accessible sizes $n = 6$--$14$ qubits?}
        \\ \emph{Basis:} We verified the \emph{order} lift, not the crossover point on any
        specific $H$. The paper's engineering claim is about competitiveness vs $S_4$ on real
        target Hamiltonians, which this replication did not test.
        \\ \emph{Next steps:} Implement $H_{\text{TFIM}} = -J\sum X_iX_{i+1} - h\sum Z_i$ for
        $n \in \{6,8,10,12\}$; run $S_1$, $S_2$, $S_4$, and Richardson MPFs at matched target
        error $\{10^{-3}, 10^{-4}, 10^{-5}\}$; log gate counts and total evolution time; find
        crossover $t^*(n, \epsilon)$ where MPF wins on gate count. Repeat for 1D
        Fermi--Hubbard via Jordan--Wigner.

  \item \textbf{How does the LCU/MPF primitive integrate with modern qubitization and QSVT ---
        is a hybrid MPF-QSVT scheme ever advantageous, or does pure QSVT strictly dominate?}
        \\ \emph{Basis:} The paper predates qubitization (2017) and QSVT (2019). The Lemma 2
        subroutine is the same block-encoding primitive but the 2012 MPF construction is not
        embedded in the modern polynomial-transform framework.
        \\ \emph{Next steps:} Build a small block-encoding of $H_{\text{TFIM}}$ ($n=4$--$6$);
        run QSVT via Chebyshev / Jacobi--Anger for $e^{-iHt}$; compare gate count and
        post-selection failure vs Lemma-2 MPF at matched target error. Attempt Richardson-in-
        degree extrapolation on QSVT polynomial approximants and check whether it lifts the
        effective order of a truncated Jacobi--Anger series.

  \item \textbf{What is the noise-robustness envelope of LCU post-selection under realistic
        gate noise: at what depolarizing / coherent error rate does
        $\Delta^2\kappa/(\kappa+1)^2$ stop being a useful predictor of ancilla-measurement
        statistics?}
        \\ \emph{Basis:} All numerics here are ideal statevector. The exact-to-$10^{-9}$
        agreement with the closed-form failure probability is purely mathematical; noise on
        controlled-$U_a/U_b$ biases the post-selected branch.
        \\ \emph{Next steps:} Reimplement Lemma 2 in a density-matrix simulator with
        depolarizing $p \in [10^{-4},10^{-2}]$ and 2-qubit gate error $\in [10^{-3},3\!\times\!10^{-2}]$;
        measure empirical failure prob vs prediction and fidelity of the success branch;
        identify the noise threshold at which post-selection no longer purifies faster than
        it accumulates error.

  \item \textbf{Can Richardson-MPF be extended cleanly to time-dependent $H(t)$, and does the
        order lift survive under Magnus / Dyson series treatment?}
        \\ \emph{Basis:} The paper treats time-independent $H$ only. Real chemistry / control
        problems are time-dependent. Whether the negative-coefficient linear combination
        $(-\tfrac{1}{3} S_2(t) + \tfrac{4}{3} S_2(t/2)^2)$ preserves its cancellation for a
        Magnus expansion is not obvious.
        \\ \emph{Next steps:} Take $H(t) = (1-s(t))H_{\text{start}} + s(t) H_{\text{final}}$
        (slow adiabatic schedule, $n=2$--$4$); implement time-ordered $S_2$-analog (midpoint
        Magnus-2); build the MPF Richardson combination; compare error scaling against a
        fine-grid Dyson-series ground truth. If order fails to lift, characterize which
        Magnus commutator terms break the cancellation and whether modified MPF coefficients
        restore it.

  \item \textbf{For the gate-count-vs-error Pareto frontier at fixed evolution time, where
        does MPF sit relative to $S_4$, $S_6$, qDRIFT, and QSVT --- and does its position
        depend on Hamiltonian sparsity / locality?}
        \\ \emph{Basis:} The paper compares MPF asymptotically only to Suzuki formulas.
        Modern practice includes qDRIFT (Campbell 2019) and QSVT. No published all-method
        Pareto sweep exists for standard benchmark Hamiltonians.
        \\ \emph{Next steps:} For each of three benchmark $H$s (TFIM $n=10$, 1D Fermi--Hubbard
        $n=6$, $H_2$/STO-3G molecular): sweep target error $10^{-2}$--$10^{-6}$; for
        $\{S_2, S_4, S_6, \text{MPF}/S_2, \text{MPF}/S_4, \text{qDRIFT}, \text{QSVT}\}$
        count exact 2-qubit gates (or controlled-Pauli exponentials); plot the Pareto
        frontier; repeat sparse-vs-dense at the same lattice size. Deliver a single figure
        answering: \emph{does 2012 MPF remain on the Pareto frontier in 2026?}
\end{enumerate}
